\documentclass[11pt]{article}
\usepackage[margin=1in]{geometry}
\usepackage[T1]{fontenc}
\usepackage{lmodern}
\usepackage{amsmath,amssymb,amsthm,mathtools}
\usepackage{hyperref}
\hypersetup{colorlinks=true,linkcolor=blue,urlcolor=blue,citecolor=blue}
\usepackage{microtype}

\newtheorem{theorem}{Theorem}
\newtheorem{proposition}{Proposition}
\newtheorem{corollary}{Corollary}
\newtheorem{remark}{Remark}
\newcommand{\Om}{\Omega}
\newcommand{\EE}{\mathcal E}
\newcommand{\RR}{\mathbb R}
\newcommand{\CC}{\mathbb C}
\newcommand{\DD}{\mathrm D}
\newcommand{\Rea}{\operatorname{Re}}
\newcommand{\Th}{\Theta_\varepsilon}

\title{E18: Exact Lyapunov Descent for the Regularized WCT Curvature Energy}
\author{Proof note for Manuscript 2}
\date{October 26, 2026}

\begin{document}
\maketitle

\begin{abstract}
We prove a precise conditional version of E18: the regularized Wave
Confinement Theory (WCT) curvature-plus-gradient energy is nonnegative and
nonincreasing along its exact negative $L^2$ gradient flow.  We compute the
full complex variational derivative, including differentiation of the
nonlinear curvature weight, and obtain an exact energy--dissipation
identity.  Periodic boundary conditions and classical solution regularity
are imposed to make every integration by parts and chain rule legitimate.
The theorem does not establish global existence of the flow or identify it
with the independently proposed physical-time WCT dynamics.
\end{abstract}

\section{Definition and assumptions}
Let $\Om=\mathbb T^d$ be a flat $d$-dimensional torus, $1\le d\le3$.
Fix $c_1,c_2\ge0$, $\varepsilon>0$, and $\alpha\in\RR$.  For a complex
field $\psi:\Om\to\CC$ set
\begin{align}
 \rho&=|\psi|^2,\\
 D(\rho)&=\rho+\varepsilon^2e^{-2\alpha\rho},\\
 R_\varepsilon(\psi)&=\frac{\overline\psi}{D(|\psi|^2)},\\
 \Th[\psi]&=-(\Delta\psi)R_\varepsilon(\psi),\\
 F(\rho)&=\frac{\rho}{D(\rho)^2}.
\end{align}
The energy is
\begin{equation}\label{eq:energy}
\boxed{\displaystyle
 \EE[\psi]=\int_\Om\left(c_1|\nabla\psi|^2
                +c_2|\Th[\psi]|^2\right)dx
 =\int_\Om\left(c_1|\nabla\psi|^2
                +c_2F(|\psi|^2)|\Delta\psi|^2\right)dx .}
\end{equation}
All variational derivatives below use the convention
\begin{equation}\label{eq:convention}
 \DD\EE(\psi)[h]
 =2\Rea\int_\Om \frac{\delta\EE}{\delta\overline\psi}
                         \overline h\,dx,
\end{equation}
where $h$ is a complex variation and the underlying tangent space is real.
In particular, $\delta\EE/\delta\overline\psi$ is one-half of the
real-$L^2$ Hilbert gradient for the inner product
$(u,v)_\RR=\Rea\int_\Om u\overline v\,dx$.

\begin{proposition}[Well-defined, nonnegative energy]\label{prop:positivity}
The denominator $D(\rho)$ is strictly positive for every $\rho\ge0$.
For every $\psi\in H^2(\Om;\CC)$, the functional in
\eqref{eq:energy} is finite and nonnegative.
\end{proposition}
\begin{proof}
As $\varepsilon>0$, one has
$D(\rho)=\rho+\varepsilon^2e^{-2\alpha\rho}>0$.
Thus $F$ is smooth and nonnegative on $[0,\infty)$.
For $d\le3$, $H^2(\Om)$ embeds continuously in $L^\infty(\Om)$;
therefore $F(|\psi|^2)$ is bounded, while
$\nabla\psi,\Delta\psi\in L^2(\Om)$.  Both terms in
\eqref{eq:energy} are finite and nonnegative because $c_1,c_2\ge0$.
\end{proof}

\section{Exact first variation}
The derivative of the curvature weight is
\begin{align}
 D'(\rho)&=1-2\alpha\varepsilon^2e^{-2\alpha\rho},\\
 F'(\rho)&=\frac{D(\rho)-2\rho D'(\rho)}{D(\rho)^3}\\
 &=\boxed{\frac{\varepsilon^2e^{-2\alpha\rho}(1+4\alpha\rho)-\rho}
                   {\bigl(\rho+\varepsilon^2e^{-2\alpha\rho}\bigr)^3}}.
 \label{eq:Fprime}
\end{align}


\begin{proposition}[Real Fr\'echet differentiability on $H^2$]
\label{prop:frechet}
For fixed $1\le d\le3$, $\varepsilon>0$, $\alpha\in\RR$, and real
$c_1,c_2$, the energy \eqref{eq:energy} is a continuously Fr\'echet
differentiable map
\[
  \EE:H^2(\mathbb T^d;\CC)_{\RR}\longrightarrow\RR.
\]
Its derivative, taken over the real scalar field, is the bounded
real-linear functional
\begin{equation}\label{eq:frechet}
\begin{split}
\DD\EE(\psi)[h]
={}&2c_1\Rea\int_\Om
       \nabla\psi\cdot\nabla\overline h\,dx\\
 &+2c_2\Rea\int_\Om
       F(|\psi|^2)\Delta\psi\,\Delta\overline h\,dx\\
 &+2c_2\Rea\int_\Om
       F'(|\psi|^2)\psi|\Delta\psi|^2\overline h\,dx .
\end{split}
\end{equation}
More precisely, for each $M<\infty$ there is $C_M<\infty$, depending
on the fixed parameters and the torus, such that whenever
$\|\psi\|_{H^2}\le M$ and $\|h\|_{H^2}\le1$,
\begin{equation}\label{eq:frechet-remainder}
 \bigl|\EE[\psi+h]-\EE[\psi]-\DD\EE(\psi)[h]\bigr|
 \le C_M\|h\|_{H^2}^2.
\end{equation}
The derivative is locally Lipschitz in the operator norm of
$\mathcal L_{\RR}(H^2,\RR)$.
\end{proposition}
\begin{proof}
Write $\delta=\|h\|_{H^2}$ and use the embedding
$H^2(\mathbb T^d)\hookrightarrow L^\infty$ for $d\le3$.
Let $q=\Delta\psi$, $r=\Delta h$, $\rho=|\psi|^2$,
$\eta=\psi\overline h+\overline\psi h$, and $\zeta=|h|^2$.
Then
\[
 |\psi+h|^2=\rho+\eta+\zeta,\qquad
 |\Delta(\psi+h)|^2
 =|q|^2+2\Rea(q\overline r)+|r|^2.
\]
On the prescribed bounded $H^2$ ball,
\[
 \|\eta\|_\infty\le C_M\delta,\quad
 \|\zeta\|_\infty\le C_M\delta^2,\quad
 \|r\|_2\le\delta,\quad
 \|q\|_2\le M.
\]
Both $\rho$ and $\rho+\eta+\zeta$ belong to a fixed compact subset
of $[0,\infty)$; positivity of $D$ implies that $F,F',F''$ are
bounded there. Taylor's theorem on the segment between these
nonnegative endpoints therefore gives
\[
 F(\rho+\eta+\zeta)=F(\rho)+F'(\rho)\eta+B,\qquad
 \|B\|_\infty\le C_M\delta^2.
\]
The linear part of
$F(\rho+\eta+\zeta)|q+r|^2-F(\rho)|q|^2$ is
\[
 2F(\rho)\Rea(q\overline r)
       +F'(\rho)\eta|q|^2.
\]
After subtracting it, the remaining terms are exactly
\[
 F(\rho)|r|^2
 +2F'(\rho)\eta\Rea(q\overline r)
 +F'(\rho)\eta|r|^2
 +B|q|^2+2B\Rea(q\overline r)+B|r|^2.
\]
Their integrals are bounded by $C_M\delta^2$ using
Cauchy--Schwarz, the displayed $L^\infty$ bounds, and $\delta\le1$.
The gradient part expands exactly as
\[
 |\nabla(\psi+h)|^2
 =|\nabla\psi|^2
   +2\Rea(\nabla\psi\cdot\nabla\overline h)
   +|\nabla h|^2;
\]
its remainder is likewise $O(\delta^2)$.
Since
$F'(\rho)\eta|q|^2
 =2\Rea(F'(\rho)\psi|q|^2\overline h)$,
this proves both \eqref{eq:frechet} and
\eqref{eq:frechet-remainder}. Each term in
\eqref{eq:frechet} is bounded by $C_M\|h\|_{H^2}$,
so the displayed formula is a real bounded linear functional.

For completeness, if $\psi,\varphi$ lie in a bounded $H^2$ ball,
the smooth scalar compositions satisfy
\[
 \|F(|\psi|^2)-F(|\varphi|^2)\|_\infty+
 \|F'(|\psi|^2)\psi-F'(|\varphi|^2)\varphi\|_\infty
 \le C_M\|\psi-\varphi\|_{H^2}.
\]
Furthermore,
\[
 \bigl\||\Delta\psi|^2-|\Delta\varphi|^2\bigr\|_1
 \le(\|\Delta\psi\|_2+\|\Delta\varphi\|_2)
       \|\Delta(\psi-\varphi)\|_2 .
\]
Insert these estimates into \eqref{eq:frechet}, using
$\|h\|_\infty\le C\|h\|_{H^2}$, to obtain
$\|\DD\EE(\psi)-\DD\EE(\varphi)\|_{\mathrm{op}}
 \le C_M\|\psi-\varphi\|_{H^2}$.
Thus $\EE$ is $C^1$ as a real Banach-space functional.
\end{proof}

\begin{proposition}[Full complex variational derivative]\label{prop:variation}
For $\psi\in H^4(\Om;\CC)$, the $L^2$ representative of the
variational derivative in convention \eqref{eq:convention} is
\begin{equation}\label{eq:G}
\boxed{\displaystyle
 G[\psi]:=\frac{\delta\EE}{\delta\overline\psi}
   =-c_1\Delta\psi
     +c_2\Delta\!\bigl(F(|\psi|^2)\Delta\psi\bigr)
     +c_2F'(|\psi|^2)\psi|\Delta\psi|^2 .}
\end{equation}
\end{proposition}
\begin{proof}
Let $h\in C^\infty(\Om;\CC)$ and $\psi_s=\psi+sh$ for real $s$.
Direct differentiation gives
\begin{align}
 \left.\frac{d}{ds}|\psi_s|^2\right|_{s=0}
 &=\psi\overline h+\overline\psi h,\\
 \left.\frac{d}{ds}|\Delta\psi_s|^2\right|_{s=0}
 &=\Delta h\,\Delta\overline\psi
   +\Delta\psi\,\Delta\overline h.
\end{align}
For the gradient part, periodic integration by parts yields
\begin{equation}\label{eq:first-gradient}
 \DD\!\left(c_1\|\nabla\psi\|_2^2\right)[h]
 =2c_1\Rea\int_\Om \nabla\psi\cdot\nabla\overline h\,dx
 =-2c_1\Rea\int_\Om\Delta\psi\,\overline h\,dx.
\end{equation}
For the curvature part the product rule gives
\begin{align}
 &\DD\!\left(c_2\int_\Om F(|\psi|^2)|\Delta\psi|^2dx\right)[h]
 \nonumber\\
 &\qquad=2c_2\Rea\int_\Om
     F(|\psi|^2)\Delta\psi\,\Delta\overline h\,dx
 +2c_2\Rea\int_\Om
     F'(|\psi|^2)\psi|\Delta\psi|^2\overline h\,dx.
 \label{eq:first-curvature}
\end{align}
The Laplacian is self-adjoint on $\mathbb T^d$, so
\begin{equation}
 \int_\Om F(|\psi|^2)\Delta\psi\,\Delta\overline h\,dx
 =\int_\Om
   \Delta\!\bigl(F(|\psi|^2)\Delta\psi\bigr)
   \overline h\,dx.
\end{equation}
Combining this with \eqref{eq:first-gradient} and
\eqref{eq:first-curvature},
\begin{equation}
 \DD\EE(\psi)[h]=2\Rea\int_\Om G[\psi]\overline h\,dx.
\end{equation}
This is \eqref{eq:G}.  The $H^4$ assumption, together with the
Sobolev algebra property in $d\le3$ and the smoothness of $F$,
ensures that every term of $G[\psi]$ lies in $L^2$.
\end{proof}

\section{Exact gradient-flow dissipation theorem}
\begin{theorem}[E18, conditional exact Lyapunov descent]\label{thm:E18}
Let $\psi\in C^1([0,T];H^4(\mathbb T^d;\CC))$ solve the specified
negative-gradient evolution
\begin{equation}\label{eq:flow}
\boxed{\displaystyle
 \partial_t\psi=-G[\psi]
 =c_1\Delta\psi
  -c_2\Delta\!\bigl(F(|\psi|^2)\Delta\psi\bigr)
  -c_2F'(|\psi|^2)\psi|\Delta\psi|^2.}
\end{equation}
Then, for every $t\in[0,T]$,
\begin{equation}\label{eq:rate}
\boxed{\displaystyle
 \frac{d}{dt}\EE[\psi(t)]
 =-2\|G[\psi(t)]\|_{L^2(\Om)}^2
 =-2\|\partial_t\psi(t)\|_{L^2(\Om)}^2
 \le0.}
\end{equation}
For any $0\le t_1\le t_2\le T$,
\begin{equation}\label{eq:integrated}
\boxed{\displaystyle
 \EE[\psi(t_2)]
 +2\int_{t_1}^{t_2}\|G[\psi(s)]\|_2^2\,ds
 =\EE[\psi(t_1)] .}
\end{equation}
In particular, $\EE[\psi(t_2)]\le\EE[\psi(t_1)]$.
\end{theorem}
\begin{proof}
Since the solution is classical in the stated Sobolev class,
the real Banach-space chain rule and Proposition~\ref{prop:variation}
give
\begin{align}
 \frac{d}{dt}\EE[\psi(t)]
 &=\DD\EE(\psi(t))[\partial_t\psi(t)]\\
 &=2\Rea\int_\Om
       G[\psi(t)]\,\overline{\partial_t\psi(t)}\,dx.
\end{align}
Inserting $\partial_t\psi=-G[\psi]$ yields
\begin{align}
 \frac{d}{dt}\EE[\psi(t)]
 &=-2\Rea\int_\Om G[\psi(t)]\overline{G[\psi(t)]}\,dx\\
 &=-2\int_\Om|G[\psi(t)]|^2dx\le0,
\end{align}
proving \eqref{eq:rate}. Integrating over $[t_1,t_2]$ proves
\eqref{eq:integrated}, and dropping the nonnegative dissipation
integral gives monotonicity.
\end{proof}

\begin{corollary}[Characterization of zero instantaneous dissipation]
Under the hypotheses of Theorem~\ref{thm:E18}, at each $t$ one has
\begin{equation}
 \frac{d}{dt}\EE[\psi(t)]=0
 \quad\Longleftrightarrow\quad
 G[\psi(t)]=0\text{ in }L^2
 \quad\Longleftrightarrow\quad
 \partial_t\psi(t)=0.
\end{equation}
Thus a nonstationary time slice has strictly negative energy derivative.
\end{corollary}
\begin{proof}
The equivalent conditions follow immediately from
\eqref{eq:rate}, the nonnegativity of a squared norm, and the
evolution equation \eqref{eq:flow}.
\end{proof}

\begin{corollary}[Finite accumulated dissipation]
If the hypotheses of Theorem~\ref{thm:E18} hold on every finite
interval $[0,T]$ and the solution exists globally for $t\ge0$, then
\begin{equation}\label{eq:infinite}
 \int_0^\infty\|G[\psi(t)]\|_2^2dt
 \le\frac12\EE[\psi(0)].
\end{equation}
\end{corollary}
\begin{proof}
By Proposition~\ref{prop:positivity} and \eqref{eq:integrated},
$2\int_0^T\|G[\psi(t)]\|_2^2dt\le\EE[\psi(0)]$ for every $T$.
Monotone convergence of the nonnegative time integral proves
\eqref{eq:infinite}.
\end{proof}

\section{Boundary terms and conventions}
\begin{remark}[Nonperiodic domains]
For a sufficiently smooth bounded domain, the first variation has,
in addition to the bulk term in \eqref{eq:G}, the boundary contribution
\begin{equation}\label{eq:boundary}
 \begin{split}
 B[h]=2\Rea\int_{\partial\Om}\bigl[&
 \bigl(c_1\partial_n\psi-c_2\partial_n(F(|\psi|^2)\Delta\psi)\bigr)
 \overline h\\
 &+c_2F(|\psi|^2)\Delta\psi\,\partial_n\overline h
 \bigr]dS.
 \end{split}
\end{equation}
The bulk $L^2$-gradient derivation applies when the admissible variations
and boundary conditions make $B[h]=0$, and the evolution is compatible
with those boundary conditions.  Imposing an unspecified boundary condition
is not enough for a fourth-order integration-by-parts argument.
\end{remark}

\begin{remark}[Factor of two]
With the convention \eqref{eq:convention}, the exact flow
$\psi_t=-\delta\EE/\delta\overline\psi$ yields the coefficient $2$
in \eqref{eq:rate}.  With the real Hilbert gradient
$\nabla_\RR\EE=2G$, the flow $\psi_t=-\nabla_\RR\EE$ instead obeys
$d\EE/dt=-\|\nabla_\RR\EE\|_2^2$.
The factor is a normalization choice, not a substantive discrepancy.
\end{remark}

\section{Scope and limitations}
Theorem~\ref{thm:E18} establishes an \emph{exact conditional}
Lyapunov identity for solutions of \emph{the specified gradient flow}.
It does not assert: (i) local or global existence for arbitrary initial
data; (ii) uniform parabolicity, since $F(0)=0$; (iii) convergence
to a unique equilibrium; (iv) whole-space spatial localization;
or (v) derivation of \eqref{eq:flow} from a distinct Hamiltonian,
wave, or mixed dissipative WCT dynamical equation.

For an arbitrary specified evolution $\psi_t=X[\psi]$, the same
chain rule gives only
\begin{equation}
 \frac{d}{dt}\EE[\psi(t)]
 =2\Rea\int_\Om G[\psi]\,\overline{X[\psi]}\,dx,
\end{equation}
and this expression does not have a fixed sign without additional
compatibility assumptions.

\noindent\textbf{Conclusion.}
The mathematically precise E18 statement is
\begin{equation*}
 \boxed{
 \psi_t=-\frac{\delta\EE}{\delta\overline\psi}
 \quad\Longrightarrow\quad
 \frac{d\EE}{dt}
 =-2\left\|\frac{\delta\EE}{\delta\overline\psi}\right\|_2^2
 \le0,
 }
\end{equation*}
for the corrected regularized energy, on the stated class of solutions.

\begin{thebibliography}{9}
\bibitem{WCTCorrected}
R.~J.~Reyes, \emph{WCT Full Equation List Corrected},
E17--E18, online repository,
\url{https://github.com/rickyjreyes/geometry_of_resonance/blob/main/WCT_FULL_EQUATION_LIST_CORRECTED.md}.

\bibitem{WCTAudit}
R.~J.~Reyes, \emph{WCT Lean Derived Audit}, E18 conditional
algebraic descent core,
\url{https://github.com/rickyjreyes/wct-lean/blob/main/WCTLean/DerivedAudit.lean}.
\end{thebibliography}
\end{document}
